\chapter{LITERATURE SURVEY} \label{c2}
\section{Overview}
Literature on attendance systems has evolved from manual practices to OTP-based, biometric, and IoT-enabled approaches. Each approach improves specific aspects but also introduces practical limitations. This chapter summarizes key approaches and identifies the gap addressed by the proposed system.

\section{Manual and Semi-Digital Attendance}
Traditional attendance methods are based on roll-call or paper registers. Semi-digital methods use spreadsheets or basic classroom apps. These methods are straightforward but lack strong verification and are highly dependent on manual checks.

\section{OTP-based Attendance}
OTP systems improve speed and reduce manual effort. Teacher shares a session code and students submit attendance. However, OTP sharing remains a major issue; if code is forwarded outside the classroom, unauthorized marking can still occur.

\section{Biometric and Face-based Attendance}
Face verification-based systems improve identity assurance and can significantly reduce impersonation. However, biometric-only systems do not independently confirm physical proximity to teacher location and may face reliability challenges under lighting and camera variations.

\section{Proximity-based Attendance Models}
Proximity-aware systems use Bluetooth, Wi-Fi, or geofencing to validate classroom presence. BLE is preferred in many mobile scenarios due to low power consumption and short-range suitability. However, proximity alone cannot confirm user identity.

\section{Recent Technical Trends}
Modern implementations increasingly combine multiple factors. As seen in recent computer-vision and deep-learning-enabled systems \cite{wang2022lip,liu2024camouflage}, practical deployments prioritize robust verification pipelines and context-aware checks. In higher education settings, student wellness and operational efficiency also influence adoption of digital workflows \cite{slavin2014medical}.

\section{Comparative Analysis}
\begin{table}[H]
\centering
\caption{Comparison of Common Attendance Approaches}
\label{c2:tab1}
\begin{tabular}{|L{3.2cm}|L{3.0cm}|L{3.0cm}|L{3.0cm}|}
\hline
\textbf{Method} & \textbf{Strength} & \textbf{Limitation} & \textbf{Proxy Risk}\\
\hline
Manual Roll Call & Simple and familiar & Time-consuming and error-prone & High\\
\hline
OTP-only & Fast and scalable & Code can be shared & Medium to High\\
\hline
Face-only & Identity-focused & No proximity assurance & Medium\\
\hline
Bluetooth-only & Presence-focused & No identity assurance & Medium\\
\hline
Multi-factor (OTP + BLE + Face) & Balanced security and usability & More integration effort & Low\\
\hline
\end{tabular}
\end{table}

\section{Research Gap}
From the surveyed methods, the key gap is the lack of integrated, classroom-practical, multi-factor validation. Most systems prioritize only one dimension: speed, proximity, or identity. The proposed Smart Attendance System addresses this by enforcing a sequence of checks before attendance insertion.

\section{Conclusion of Literature Survey}
The survey indicates that multi-factor attendance is better suited for real classroom conditions than single-factor systems. Therefore, combining OTP session control, BLE proximity, and face verification forms a practical and secure attendance framework.

